\subsection{Actual repairs supported in high weights}

Throughout this subsection $R$ is $\mathbb Z$ or
$\mathbb Z_{(p)}$, and $S_R=\Gamma_R(T)$.
For $n\geq1$, write $D(n)$ for the Dold--Kan image of
the identity chain complex $R\to R$ in homological
degrees $n,n-1$.  This simplicial module is projective:
under normalization, a chain map from that disk to a
complex $C$ is specified by an arbitrary element of
$C_n$, and evaluation at $n$ is exact.  It is also
contractible, and $D(n)_0=0$ for $n\geq2$.
The term ``free PD algebra on a graded $S_R$-module''
below means the relative graded algebra with divided
powers on its even generators and exterior powers on
its odd generators.  Its universal property imposes
the given $S_R$-linear relations on the generators;
we do not assume that those modules are free over $S_R$.

\begin{proposition}[AT13: Proved; coefficient-tail repair with an explicit horizontal retraction]
\label{AT13}
Let $X\to A,B$ be actual strict simplicial PD-DGA maps,
where $A,B$ are admissible over $R$.  Let $q\geq1$,
$b\geq0$, $N\geq1$, $M=D(n)$ with $n\geq1$, and put
$L=I_N\langle b\rangle$ and $W=N+b$.
Suppose $W\geq q$ is above the uniform vertical
cutoffs of $A,B$ through form degree $q$.
For every specified simplicial $S_R$-linear map
\[
                 z:L\otimes_R M\longrightarrow Z^q(X)
 \tag{AT.50}
\]
there is an actual strict PD extension $X\hookrightarrow Y$
with maps $Y\to A,B$ extending the specified maps.
Every element in the image of \textup{(AT.50)} becomes
a vertical boundary.  In all form degrees the extension
is an isomorphism in weights less than $W$.
Every horizontal form-weight row of $Y$ deformation
retracts onto the corresponding row of $X$.

For $q\geq2$, suppose in addition that $X,A,B$ have
torsionfree $R$-terms and that the coefficient map gives
$H^0(X_n)=S_R$ for every $n$.  Set
$Y^{\mathrm{tf}}=\operatorname{im}(Y\to Y\otimes_R\mathbb Q)$.
The maps and horizontal retractions descend to this
quotient, its terms are flat over $R$, and
$H^0(Y^{\mathrm{tf}}_n)=S_R$ in every weight and simplex.
This last assertion does not assert vanishing of its
positive vertical cohomology.
\end{proposition}

\begin{proof}
Choose a graded free resolution $Q_\bullet\to L$
whose generators all have weights at least $W$.
The weightwise construction of AT11 gives such a
resolution, with finite free $R$-weight groups and
free weight groups in every kernel.  Put $K_0=L$
and $K_i=\ker(Q_{i-1}\to K_{i-1})$ for $i\geq1$.
For $q=1$ put $V^0=L$.
For $q\geq2$, define the graded module of new
generators by
\[
 V^{q-1}=Q_0,\quad V^{q-2}=Q_1,\quad\ldots,
 \quad V^1=Q_{q-2},\quad V^0=K_{q-1}.
 \tag{AT.51}
\]
The repeated notation in \textup{(AT.51)} has its
literal indicated range: when $q=2$ it means
$V^1=Q_0$, $V^0=K_1$.
The maps between these terms are the resolution
differentials, ending with the inclusion of $K_{q-1}$.

As a simplicial graded PD algebra, put
\[
 Y=X\otimes_{S_R}\operatorname{FreePD}_{S_R}(V\otimes_R M).
 \tag{AT.52}
\]
Give the top new generators differential $z\epsilon$,
where $\epsilon:Q_0\to L$, and give the remaining
new generators the differentials in \textup{(AT.51)}.
For $q=1$ the formula is simply $d(v)=z(v)$.
Extend by the derivation and PD differentiation rules.
These assignments are $S_R$-linear, preserve all
generator relations, and square to zero, because
$z$ is closed and the displayed resolution is a
complex.  They therefore define a differential on
the free PD presentation.  The construction is
simplicial because \textup{(AT.50)} and every linear
assignment are simplicial.  This defines the
claimed strict PD extension, including its differential.
The augmentation of the free factor gives a graded
simplicial retraction onto $X$, so the inclusion is
injective.  That retraction is used only on horizontal
rows; it need not commute with the vertical differential.
All generators have weights at least $W$, which proves
the claimed isomorphism in lower weights and preserves
$q\leq w$ when it holds for $X$.
Surjectivity of $Q_0\otimes_R M_a\to L\otimes_R M_a$
in every simplex $a$ proves the assertion about boundaries.

We construct the two target maps, rather than assume
coherent primitives.  For either target $A$, consider
the total Hom complex from $Q_\bullet\otimes_R M$
to its cochain complex of simplicial graded modules.
For a graded free term $Q_i$ this Hom functor is a
product of functors
$\operatorname{Hom}_{s\operatorname{Mod}_R}(M,A^w)$,
where $w\geq W$.  Each is exact in its target because
$M=D(n)$ is projective.  Thus its vertical cohomology
is the corresponding Hom into $H^j(A)$.
For $1\leq j\leq q$ it is zero by the specified
uniform vertical cutoff.  For $j=0$ the coefficient
map identifies the target with the constant simplicial
module $S_R^w$, and
\[
 \operatorname{Hom}_{s\operatorname{Mod}_R}(D(n),S_R^w)=0
                      \qquad(n\geq1),
 \tag{AT.53}
\]
since the normalized target has zero term in degree $n$.
The first-quadrant Hom bicomplex consequently has zero
cohomology in total degree $q$: all vertical cohomology
terms on that finite diagonal vanish.
The composite of \textup{(AT.50)} with $X\to A$
therefore has an actual nullhomotopy on this resolution.
After the usual successive sign adjustment, its
components are simplicial $S_R$-linear maps
\[
 h_i:Q_i\otimes_R M\longrightarrow A^{q-1-i},
 \quad d h_0=fz\epsilon,\quad
 d h_i=h_{i-1}\partial,\quad h_{q-1}\partial_q=0.
 \tag{AT.54}
\]
The final equation makes the last map factor through
$K_{q-1}\otimes_R M$; tensoring here is exact since
each $M_a$ is free over $R$.
For $q=1$ the same argument factors $h_0$ through
$L\otimes_R M$.
These are precisely the generator values required
for a strict map \textup{(AT.52)} to $A$.
Repeating the construction for $B$ gives the other
actual map.

Let $h_i^M:M_a\to M_{a+1}$ be an additive simplicial
homotopy from the identity to zero.  Define on
\textup{(AT.52)} the homotopy operators which are
$s_i$ on the $X$ factor and
$\operatorname{FreePD}(1_V\otimes h_i^M)$ on the
free factor.  The coefficients are fixed, so these
maps descend through the tensor product over $S_R$.
The simplicial homotopy identities are equalities
of composites and therefore are preserved by the
free PD functor, including with the stated nonfree
generator modules.  Their endpoints are the identity
of $Y$ and its graded augmentation followed by the
inclusion of $X$.
They preserve form degree and weight and hence give
the asserted deformation retraction of every row.
No assertion that these operators commute with the
vertical differential is used.

For the torsionfree assertion, the $R$-torsion in
$Y$ is an ideal preserved by the differential,
simplicial maps, and divided powers: if $a$ is
annihilated by $r\ne0$, then $\gamma_m(a)$ is
annihilated by $r^m$.
Its quotient is the displayed image $Y^{\mathrm{tf}}$.
Both target maps kill that ideal.  All horizontal
homotopy operators preserve it, and the retraction
lands in the torsionfree module $X$, so the entire
row retraction descends.  Torsionfree modules over
the PID $R$ are flat.

Finally rationalize the coefficient resolution.
The ideal $L_{\mathbb Q}$ is a free graded
$\mathbb Q[t]$-module on $T^{[N]}$ with its shift.
The resolution splits into this module and identity
pairs: successively each kernel is graded projective,
so each displayed surjection splits in graded modules.
The splits are in the constant coefficient modules
and hence commute with the simplicial parameter $M$.
The rational algebra \textup{(AT.52)} is consequently
$X_{\mathbb Q}$ with generators in form degree
$q-1$ and free contractible disks in consecutive
form degrees.  A disk contracts over $\mathbb Q$
in every positive arity, by the even-to-odd formula
of C21 or by division by its nonzero odd-to-even
arity coefficient.  Its contraction is linear over
the remaining algebra.  Since $q-1\geq1$, the
remaining generators do not change form degree zero
or its cycles.  Hence inclusion identifies
$H^0(Y_{a,\mathbb Q})=S_{\mathbb Q}$ in each simplex.

In weights below $W$ the original $X$ was unchanged.
In a weight at least $W$, a closed element of
$(Y^{\mathrm{tf}})_a^0$ is, by the injective map to
the rational algebra, an actual rational coefficient:
there are no form-degree-minus-one boundaries.
Its image in $A_a$ is closed, and the vertical
coefficient comparison there says it is an element
of $S_R$ in that weight.  Rationalizing this equality
and using the coefficient-preserving target map
shows that the original rational coefficient was
already in $S_R$.
Injectivity of $Y^{\mathrm{tf}}\to Y_{\mathbb Q}$
now gives the asserted equality of $H^0$ integrally.
\end{proof}

\begin{theorem}[AT14: Proved; a two-cell repair preserving the rational coefficient resolution]
\label{AT14}
Let $X\to A,B$ be actual strict simplicial PD-DGA maps
over $R$, with torsionfree terms, and suppose $A,B$
are admissible.  Suppose the coefficient map gives
$H^0(X_a)=S_R$ and gives a rational quasi-isomorphism
$S_{\mathbb Q}\to X_{a,\mathbb Q}$ in every simplex
and weight.  Fix $q\geq2$, $n\geq1$, and a weight
$w\geq q$ above both target vertical cutoffs through
form degree $q$.
For every simplicial linear closed map
\[
                    z:D(n)\longrightarrow Z^q(X^{w,*}),
 \tag{AT.55}
\]
there is a strict PD extension $X\hookrightarrow Y$
with actual extending maps $Y\to A,B$ such that:
\begin{enumerate}
\item the specified cycle image becomes a vertical boundary;
\item all horizontal rows deformation retract to those of $X$,
      and all weights less than $w$ are unchanged;
\item the terms are torsionfree, and are free over $R$
      when those of $X$ are free;
\item $X_{a,\mathbb Q}\to Y_{a,\mathbb Q}$ is a
      quasi-isomorphism, and $H^0(Y_a)=S_R$ in every
      simplex and weight;
\item when $n\geq2$, the entire simplex-zero algebra is
      unchanged: $Y_0=X_0$.
\item for $q\geq3$, inclusion identifies $H^1(X_a)$
      with $H^1(Y_a)$ in every weight and simplex.
\end{enumerate}
In particular every positive vertical cohomology
group of $Y$ is $R$-torsion.  The construction does
not assert that those groups vanish.
\end{theorem}

\begin{proof}
Put $M=D(n)$.  Projectivity of $M$ makes
$\operatorname{Hom}_{s\operatorname{Mod}_{\mathbb Q}}
(M_{\mathbb Q},X_{\mathbb Q}^{w,*})$
have the cohomology of Hom into the constant
coefficient module.  It is zero in positive form
degrees, and in degree zero by \textup{(AT.53)}.
Thus $z$ has a rational simplicial primitive.
A map from $D(n)$ is determined by the image of
one normalized generator in degree $n$.
That image is an element of $X\otimes_R\mathbb Q$,
so one nonzero integer $m$ clears its denominators.
All its faces and degeneracies then have the same
denominator bound.  Torsionfreeness of $X$ turns
the resulting rational equation into an integral
one.  We obtain an actual simplicial
linear map $b:M\to X^{w,q-1}$ with
\[
                              db=mz.
 \tag{AT.56}
\]

Adjoin free PD generator copies $u(M)$ and $v(M)$,
both of weight $w$, in form degrees $q-1$ and $q-2$.
Define
\[
                     du=z,\qquad dv=mu-b.
 \tag{AT.57}
\]
Equation \textup{(AT.56)} proves $d^2v=0$;
$z$ being closed proves $d^2u=0$.
The free strict PD presentation therefore gives
an actual differential algebra $Y$ with these maps,
and the formulas are simplicial.
Its underlying algebra is
\[
 X\otimes_R
 \operatorname{FreePD}_R\bigl(M[q-1]\oplus M[q-2]\bigr),
 \tag{AT.58}
\]
where the bracket denotes placement in form degree
and both copies have the stated weight.
This is a free extension by degreewise free parameter
modules, so its terms have the stated freeness or
torsionfreeness, and its lower weights are unchanged.
The contraction of $M$ gives the explicit horizontal
retraction of \textup{(AT.58)}, by the same operator
formula as in AT13.

For a target $A$, projectivity of $M$ and its
vertical comparison through degree $q$ give a
simplicial primitive $u_A:M\to A^{w,q-1}$ of $fz$.
The map $m u_A-fb$ is closed in form degree $q-1$.
Since $q-1\geq1$, the same exact Hom calculation
gives a simplicial primitive
$v_A:M\to A^{w,q-2}$.
These two maps satisfy exactly \textup{(AT.57)},
so they define the actual extending strict PD map
$Y\to A$.  The identical procedure gives $Y\to B$.
Only the stated uniform target bounds have been used.

Over $\mathbb Q$ make the explicit invertible
change of generator $u'=u-b/m$, keeping $v$ and
$X$ fixed.  Its inverse sends $u$ to $u'+b/m$;
the free PD universal property defines both maps,
including when $u$ has even degree.
Then
\[
                         du'=0,\qquad dv=mu'.
 \tag{AT.59}
\]
Rescaling the second generator by the unit $m$
identifies this with an identity disk on $M$.
In each simplex choose a basis of $M_a$.
The resulting tensor product of one-generator
disks has a contraction on every positive arity:
for an even first generator its PD differential
is an identity pair, and for an odd first generator
the coefficient is a positive integer and is a
unit over $\mathbb Q$.
Tensor these contractions, including their ordinary
Koszul signs, to obtain an $X_{a,\mathbb Q}$-linear
deformation retraction onto $X_{a,\mathbb Q}$.
Thus the actual coefficient-preserving inclusion
is a rational quasi-isomorphism in every weight.

We have proved $H^0(Y_{a,\mathbb Q})=S_{\mathbb Q}$.
The terms of $Y$ are already torsionfree.
The weight cutoff and the actual map to $A$ now
give $H^0(Y_a)=S_R$ by the same rational-coefficient
integrality argument at the end of AT13.
Since all positive rational cohomology vanishes,
flat scalar extension to $\mathbb Q$ makes each
positive integral cohomology group zero; every
element of those groups is therefore $R$-torsion.
Finally $M_0=0$ for $n\geq2$, so neither generator
copy adds anything at simplex zero in that case.
For $q\geq4$, both new generator families have form
degree at least two, and the form-zero and form-one
terms, including their outgoing differential in form
one, are unchanged.  For $q=3$, only the $v$ family
has form degree one.  A form-one element is a sum
of an old element and a linear combination of the
$v$ generators with coefficients in $X^0$.
In its differential, the coefficient of each
independent new $u$ generator is $m$ times the
corresponding coefficient of $v$.
Torsionfreeness makes all these coefficients zero
when the element is closed.  Thus the degree-one
cycles are precisely the old ones; the boundaries
are unchanged because form degree zero is unchanged.
This proves the final assertion and completes the
construction.
\end{proof}

For $q=2$ there is also a precise calculation in the
first affected weight.  Suppose $X^{0,*}=R[0]$.
In weight $w$ no product of two new generators can
occur, so the quotient of cochain complexes is
$M\otimes_R(Rv\xrightarrow{m}Ru)$ in form degrees
zero and one.  Its only cohomology is $M/mM$ in
degree one.  The long exact sequence therefore is
the following exact sequence of simplicial modules:
\[
\begin{split}
0\longrightarrow H^1(X^{w,*})\longrightarrow H^1(Y^{w,*})
 &\longrightarrow M/mM
 \xrightarrow{\ [a]\mapsto[az]\ }H^2(X^{w,*})\\
 &\longrightarrow H^2(Y^{w,*})\longrightarrow0 .
\end{split}
 \tag{AT.60}
\]
The displayed connecting map is computed by applying
$d$ to the lift $a u$, and is well-defined because
$m z=db$.  Thus the exact new group in degree one
is the kernel of this class map.  In particular,
choosing $m$ to be the order of one class in one
simplex does not prove that the map is injective
in all simplices.  For $M=D(1)$ and $X_0=S_R$,
the bottom cycle is zero, so in simplex zero the
entire $M_0/mM_0$ is in that kernel.

The two formulas in \textup{(AT.57)} also make the
remaining low-degree issue explicit.  When $q=2$
they put $v$ in form degree zero.  If $n=1$, its
parameter has a nonzero simplex-zero term, so
simplex zero can acquire an integral disk with
nonunit differential $m$.  Formula \textup{(AT.59)}
proves that its cohomology is rationally zero;
it supplies no integral contraction when $m$ is
a nonunit.  The theorem establishes this individual
repair and its exact bounds, without asserting that
an iteration eliminates every integral vertical class.
